\documentclass[11pt]{article}
\usepackage[margin=1in]{geometry}
\usepackage{amsmath,amssymb,amsthm}
\usepackage{booktabs}
\usepackage{hyperref}
\newtheorem{theorem}{Theorem}
\newtheorem{lemma}{Lemma}
\newtheorem{proposition}{Proposition}
\newtheorem{corollary}{Corollary}
\title{Transport Bounds for Fairness Statistics: An Accounting Audit
Across Two Monitoring Programmes and the Documented Failure of
Threshold-Crossing Forecasting}
\author{Ranadhir Ghosh}
\date{}
\begin{document}
\maketitle

\begin{abstract}
When a deployed classifier is monitored for fairness, the quantities
under audit are fairness statistics --- the equalised-odds (EO) and
demographic-parity (DP) gaps --- yet the transport theory available
to monitoring has priced a different object, the false-alarm rate
(FAR) of an alarm rule. This paper closes that gap. A single
stratum-wise lemma --- the event-probability bound that the FAR
theory already rests on, applied conditionally --- yields transport
bounds for the gap statistics themselves: the DP transport between
two populations is at most the sum of the two group-conditional
score-law total variations, and the EO transport is at most the
maximum of the label-stratified transport sums, in exactly the
max-form in which the pipelines compute the EO gap. The forms are
tight as functions of the stratum variations alone, and a prevalence
decomposition reduces the DP bound to FAR-type quantities plus named,
measurable terms. We then transport the field's accounting protocol
--- per-fold comparison of realised transport against distributional
transport, at a fixed criterion --- from the FAR to the gap
functionals and run it on two monitoring programmes: US
income/employment/coverage classifiers across six states, and UK
police stop-and-search risk scores across nine forces. The recorded
FAR-accounting verdicts whose fold artifacts survive (three of five)
are reproduced under the gap bound; the family-1 gap criterion is
met on the folds that survive there (median transport quantity
0.110, 6/6 folds); the two verdicts whose folds were lost in an
archive migration are attributed to the FAR accounting only, and
said so. On the second programme the gap accounting is correct but
loose --- the gaps there are small and barely move --- and the FAR
accounting remains the binding evidence. The paper's second half is
the negative record this audit was built to serve: a five-year
programme to forecast fairness-breach crossings failed its
pre-registered falsifier on family 1 twice (the second time under a
scale-adapted design built to exclude the first failure's
explanation), and twice failed to produce a testable cell on family
2, for two distinct reasons that we prove are theorems about the
monitoring design --- crossing rarity and cut degeneracy --- rather
than implementation accidents. A fitted-model analysis then shows
the failure is consistent with, but not explained by, a
history-value bound: at family 1's measured regime the bound
permits the record. The forecasting claim is retired; what survives
is a bound, an audit, and a documented boundary.
\end{abstract}

\section{Introduction}

Fairness monitoring of a deployed classifier produces, at each
monitoring step, a small set of audited statistics --- in this
paper, the equalised-odds and demographic-parity gaps of a frozen
scoring rule on labelled batches. A natural theoretical question
for any such programme is accounting: when the audited statistic
moves between two populations (two jurisdictions, two periods, a
training pool and a deployment site), can the movement be priced by
how far the underlying distributions have moved? The one transport statement available for this purpose prices the wrong object: it bounds the change in a threshold rule's \emph{false-alarm rate} between states by the total-variation (TV) distance of the states' null score distributions. That statement is correct and classical --- it is the definitional single-event TV bound, scoped to FAR transport in this programme's own upgrade cycle, with its provenance stated exactly in Section~\ref{sec:bound} and no novelty claimed for it anywhere in this programme. But an auditor who wants
to know whether an observed \emph{gap} movement is within what
distributional transport allows finds no bound for the gap in the
text: the theorem's functional and the audited functional are
different objects.

This paper's first contribution closes that mismatch, and does so
without new machinery. Section~\ref{sec:bound} shows that the FAR
bound's own single-event argument, applied stratum-wise, composes
into transport bounds for the DP and EO gaps in exactly the forms
the monitoring pipelines compute --- including the EO gap's
max-of-components form, which the bound matches with a maximum of
transport sums. The bounds are tight in the stratum TVs, and the DP
bound admits a reduction to FAR-type (label-0) transports plus
prevalence terms, so the two theories are one theory at two
granularities.

The second contribution is an audit. An accounting protocol ---
per-fold, realised transport against distributional transport,
judged at a criterion fixed before outcomes --- existed for the
FAR. Section~\ref{sec:accounting} transports it to the gap
functionals and runs both accountings on the two programmes this
work has monitored: family 1, classifiers over US American
Community Survey tasks in six states; family 2, risk scores over UK
police stop-and-search records in nine forces. Three of the five
recorded FAR verdicts reproduce under the gap bound on their own
folds; two cannot be re-adjudicated at all, because their folds
were lost, and the paper says which and why rather than letting the
stronger theorem silently absorb them.

The third contribution is a negative result, reported as one. The
programmes existed to test whether fairness-breach crossings can be
\emph{forecast} --- alarmed before they occur, at a controlled
false-alarm rate, with a lead time beating trivial extrapolation.
That claim was pre-registered with a binding falsifier, and it
failed: on family 1 at the original absolute threshold, and again
under a scale-adapted threshold design built after the first
failure to exclude its most plausible excuse. On family 2 the test
could not even be run twice, for two distinct structural reasons
--- too few crossings exist to score, and the deployed scores
degenerate below the decision cut --- which
Section~\ref{sec:vacuity} states as theorems about monitoring
designs rather than as bad luck. Section~\ref{sec:boundary} asks
whether a history-value bound explains the forecasting failure and
answers honestly: at the measured regime it does not; it permits
the record. The claim is retired in Section~\ref{sec:forecasting}.

Nothing in the theory is novel as mathematics, and the paper does
not pretend otherwise: the stratum lemma is the classical
single-event TV bound. The contribution is the composition, the
audit at two granularities across two real monitoring programmes,
and a failure record complete enough to be useful to whoever tries
next.

\section{Monitored statistics and monitoring programmes}
\label{sec:objects}

A population is a law over records $(S, A, Y)$: $S = s(X) \in [0,1]$
the output of one fixed deployed scoring rule $s$ (the same
function for every population --- the monitoring scenario is a
frozen model under drift), $A \in \{0,1\}$ a binarised protected
attribute, $Y \in \{0,1\}$ the label. The decision is
$\hat Y = \mathbf 1\{S \ge c\}$ at the pipelines' decision cut
$c = 0.5$. For a stratum $C \in \{A{=}a\} \cup \{(Y{=}y, A{=}a)\}$,
write $R_p(C) = \Pr_p[\hat Y = 1 \mid C]$ and
$\mu_p^C = \mathrm{Law}_p(S \mid C)$. The monitored statistics, in
the exact form both programmes' pipelines compute them, are
\[
\mathrm{DP}_p = \bigl| R_p(A{=}1) - R_p(A{=}0) \bigr|,
\qquad
\mathrm{EO}_p = \max\Bigl(
\bigl| R_p(1,1) - R_p(1,0) \bigr|,\;
\bigl| R_p(0,1) - R_p(0,0) \bigr| \Bigr),
\]
where $R_p(y,a) = R_p(Y{=}y, A{=}a)$: the EO gap is the
\emph{maximum} of the true-positive-rate and false-positive-rate
differences, not either component alone, and a component whose
stratum is empty in a sample is dropped by the pipelines (the
population reading: the component is undefined; see the scope
remark in Section~\ref{sec:bound}).

\paragraph{Family 1.} Three ACS tasks \cite{acs} (income,
employment, public coverage) $\times$ six states (CA, FL, NY, PA,
TX, WA) $\times$ years 2014--2019. One frozen logistic regression per task, trained
on pooled 2014 data, never retrained; monitoring steps are years.
Family 1's original build is also the source of the forecasting
programme of Section~\ref{sec:forecasting}.

\paragraph{Family 2.} UK police stop-and-search records
\cite{policeuk} (data.police.uk), calendar 2024 (confirmatory) and
2025 (a second, untouched year used by the adapted design of
Section~\ref{sec:forecasting}). An inclusion rule frozen before
outcomes admitted 9 of 13 candidate forces for 2024 (71{,}347 usable
records) and 8 for 2025 (65{,}826). Two model classes are monitored:
logistic regression (LR) and gradient boosting (GB), each trained
once on pooled base-month data and frozen; monitoring steps are
months; the protected attribute is the recorded ethnicity
binarisation of the source data. Family 2 was constructed as an
independent second family for the accounting audit, with no overlap
with family 1's design decisions beyond the shared protocol.

\paragraph{The forecasters under audit.} The forecasting programme (Section~\ref{sec:forecasting}) evaluates candidate alarm rules against trivial baselines at a matched false-alarm rate. All forecasters act on instances (trajectory, step), whose label is whether the trajectory's gap crosses the threshold within the horizon $H$, scored from batch-observable step features: current and lagged values of the monitored gap, one-step changes, history slopes, the score-distribution gap between groups, group-share and prevalence shifts and their changes, and per-metric drift budgets. The control forecaster, CTRL --- the programme's original proposed forecaster, unchanged --- is a logistic regression on those features, Platt-calibrated (sigmoid calibration fitted by 3-fold cross-validation on the training trajectories), alarming when the calibrated score exceeds a single global threshold per (cell, FAR target), the threshold chosen on the training trajectories by the pick rule: the threshold whose achieved training FAR is closest to the target, ties toward the smaller FAR. The primary candidate, TRA-Q, uses the \emph{same fitted model as CTRL in the same fold} --- identical instance scores; only the alarm rule differs. TRA-Q alarms at step $t$ of a trajectory in target state $s$ iff $p_t \ge \theta_{s,t}$, where $\theta_{s,t}$ is the weighted $(1-q)$ quantile of a two-component pool: the global pool $G$ of all training-fold instance scores, each point weighted $(1-\lambda_t)/|G|$, and the trajectory's own pool $O_t$ of its scores at steps from the series start through $t$ inclusive, each point weighted $\lambda_t/|O_t|$, with $\lambda_t = |O_t| / (|O_t| + \kappa)$ and $\kappa = 3$. The level $q$ is calibrated per (cell, FAR target, fold) by simulating the full rule on the training trajectories, choosing the $q$ minimising $|\text{achieved training FAR} - \text{target}|$, ties toward the smaller FAR, and is then frozen for the test state. The falsifier's operating point is FAR target $0.10$; targets $0.05$ and $0.20$ are computed and reported without headline force, and the headline cells use horizon $H = 2$ steps. Baselines, each calibrated on the same training trajectories to the same FAR targets: last-value persistence, linear extrapolation of the gap trajectory, and a CUSUM detector (slack $0.02$, the original threshold grid); FairCanary-style point and interval detectors and a climatology reference were also computed. Full specifications are frozen in the programme's pre-registrations, deposited with the artifacts (Section~\ref{sec:repro}).

\section{The transport bound for the gap statistics}
\label{sec:bound}

\begin{lemma}[Stratum lemma]
\label{lem:stratum}
For every stratum $C$ with $\Pr_p(C), \Pr_q(C) > 0$,
$|R_p(C) - R_q(C)| \le \mathrm{TV}(\mu_p^C, \mu_q^C)$.
\end{lemma}

\begin{proof}
$R_p(C) = \mu_p^C([c,1])$ and $R_q(C) = \mu_q^C([c,1])$ --- the
same event under both laws, because $s$ and $c$ are state-blind ---
and a single event's probability difference is at most the total
variation, by its variational definition.
\end{proof}

The FAR transport statement is this lemma with $C$ the null-step population and $\mu$ a forecaster's step-score law: a false-alarm rate is one event's probability under one conditional score law. Its provenance, stated exactly because later sections lean on it: the inequality is the definitional single-event TV bound and is classical; its scoping to the false-alarm-rate transport of a fairness-monitoring alarm is this programme's own, fixed in the upgrade cycle's pre-registration (2026-10-29), for which a concept-level prior-art check --- recorded in that pre-registration, and described there as concept-level rather than exhaustive --- found no published statement of a bound or impossibility for a monitoring alarm's FAR transport across sites in terms of the null score distribution, and none in the fairness-monitoring setting. The classical neighbouring form is the domain-adaptation bound of Ben-David et al.~\cite{bendavid2010}, which prices a target \emph{error} by a source error plus a divergence plus a shared-hypothesis term --- a different functional, as here. No novelty is claimed for the inequality itself; the gap bounds below are the same argument composed over the strata that define the gaps.

\begin{theorem}[DP transport]
\label{thm:dp}
$|\mathrm{DP}_p - \mathrm{DP}_q| \le D_{\mathrm{DP}}(p,q) :=
\mathrm{TV}(\mu_p^{A=1}, \mu_q^{A=1}) +
\mathrm{TV}(\mu_p^{A=0}, \mu_q^{A=0})$.
\end{theorem}

\begin{proof}
By the reverse triangle inequality and then the triangle
inequality,
$|\mathrm{DP}_p - \mathrm{DP}_q|
\le |(R_p(A{=}1) - R_q(A{=}1)) - (R_p(A{=}0) - R_q(A{=}0))|
\le \sum_a |R_p(A{=}a) - R_q(A{=}a)|$,
and Lemma~\ref{lem:stratum} bounds each term.
\end{proof}

\begin{theorem}[EO transport]
\label{thm:eo}
Let $D_{\mathrm{TPR}} = \sum_a \mathrm{TV}(\mu_p^{1,a}, \mu_q^{1,a})$
and $D_{\mathrm{FPR}} = \sum_a \mathrm{TV}(\mu_p^{0,a}, \mu_q^{0,a})$.
Then the component transports satisfy
$|\mathrm{TPRgap}_p - \mathrm{TPRgap}_q| \le D_{\mathrm{TPR}}$ and
$|\mathrm{FPRgap}_p - \mathrm{FPRgap}_q| \le D_{\mathrm{FPR}}$, and
\[
|\mathrm{EO}_p - \mathrm{EO}_q| \le D_{\mathrm{EO}}(p,q) :=
\max\bigl(D_{\mathrm{TPR}},\, D_{\mathrm{FPR}}\bigr).
\]
\end{theorem}

\begin{proof}
Each component is a two-group difference of stratum rates, so
Theorem~\ref{thm:dp}'s argument within the stratum family gives
the component bounds. For the max form,
$|\max(u,v) - \max(u',v')| \le \max(|u-u'|, |v-v'|)$ for
nonnegative $u,v,u',v'$ (if $\max(u,v) \ge \max(u',v')$, say $\max(u,v) = u$, then $\max(u,v) - \max(u',v') \le u - u' \le |u-u'|$; the other case is symmetric), applied to
the two components.
\end{proof}

\paragraph{Scope remark.} The statements assume the strata
involved have positive mass under both populations. The component
bounds hold family by family whenever that family's strata are
positive in both; where a stratum has mass zero in one population,
the corresponding gap component is undefined in the pipelines too
--- statistic and bound lose the same component at the same place.
Section~\ref{sec:vacuity} meets the empirical extreme of this
condition.

\begin{proposition}[Reduction to FAR-type quantities]
\label{prop:reduction}
With prevalences $\pi_p(a) = \Pr_p[Y{=}1 \mid A{=}a]$, the law of
total probability gives
$R_p(A{=}a) = \pi_p(a) R_p(1,a) + (1-\pi_p(a)) R_p(0,a)$, and hence
$|\mathrm{DP}_p - \mathrm{DP}_q|$ is at most
$\sum_a \bigl\{ \pi_p(a)\,\mathrm{TV}(\mu_p^{1,a}, \mu_q^{1,a})
+ (1-\pi_p(a))\,\mathrm{TV}(\mu_p^{0,a}, \mu_q^{0,a})
+ |\pi_p(a) - \pi_q(a)| \cdot |R_q(1,a) - R_q(0,a)| \bigr\}$:
the DP transport is bounded by label-stratified transports --- the
label-0 terms being exactly the FAR-type objects of
Lemma~\ref{lem:stratum} --- plus named prevalence terms, every one measurable from the same records as the gaps.
\end{proposition}

\begin{proof}
Fix $a$, and abbreviate $R_p = R_p(A{=}a)$, $\pi_p = \pi_p(a)$,
$R_p(y) = R_p(y,a)$, and likewise for $q$. By the law of total
probability over $Y$, $R_p = \pi_p R_p(1) + (1-\pi_p) R_p(0)$ and
$R_q = \pi_q R_q(1) + (1-\pi_q) R_q(0)$. Adding and subtracting
$\pi_p R_q(1) + (1-\pi_p) R_q(0)$ gives the identity
\[
R_p - R_q = \pi_p \bigl[R_p(1) - R_q(1)\bigr]
+ (1-\pi_p) \bigl[R_p(0) - R_q(0)\bigr]
+ \bigl[\pi_p - \pi_q\bigr] \cdot \bigl[R_q(1) - R_q(0)\bigr],
\]
exactly --- an identity in the three displayed differences, not an
inequality. Taking absolute values, using the triangle inequality
term by term and the nonnegativity of the weights $\pi_p$ and
$1-\pi_p$, and applying Lemma~\ref{lem:stratum} to the two rate
differences, $|R_p(y) - R_q(y)| \le
\mathrm{TV}(\mu_p^{y,a}, \mu_q^{y,a})$, gives the per-group bound;
summing over $a \in \{0,1\}$ and applying the triangle inequality
once more to $\mathrm{DP}_p - \mathrm{DP}_q = [R_p(A{=}1) -
R_q(A{=}1)] - [R_p(A{=}0) - R_q(A{=}0)]$ yields the displayed bound.
The label-0 terms are the FAR-type objects: $R(0,a)$ is a label-0
conditional alarm rate, priced by the same single-event argument
as the FAR itself.
\end{proof}

\begin{corollary}[Covariate form]
Both theorems hold with every score law $\mu$ replaced by the corresponding covariate law $\nu = \mathrm{Law}(X \mid C)$, in the same sum and max forms.
\end{corollary}

\begin{proof}
Each score law is the pushforward of the corresponding covariate
law through the common map $s$: $\mu^C = \nu^C \circ s^{-1}$.
Total variation contracts under a common channel (data
processing): for every measurable $E$,
$|\mu^C(E) - \mu'^C(E)| = |\nu^C(s^{-1}E) - \nu'^C(s^{-1}E)|
\le \mathrm{TV}(\nu^C, \nu'^C)$, and taking the supremum over
$E$ on the left gives $\mathrm{TV}(\mu^C, \mu'^C) \le
\mathrm{TV}(\nu^C, \nu'^C)$. Substituting the covariate TVs for
the score TVs throughout Theorems~\ref{thm:dp} and~\ref{thm:eo}
--- the right-hand sides weakly increase, the forms are unchanged
--- gives the claim.
\end{proof}

\begin{proposition}[Tightness of form]
The sum and max forms are best possible as functions of the stratum
TVs alone: for any stratum-TV profile with $D \le 1$ there are
population pairs attaining equality (two-point score laws per
stratum, the unaffected strata held identical). Where $D > 1$ the bounds are vacuous and no tightness claim is needed.
\end{proposition}

\begin{proof}
Hold every stratum not under construction identical across $p$
and $q$; its TV and its rate difference are then $0$, so the
constructions compose stratum by stratum. \emph{Single stratum.}
Fix $d \in [0,1]$ and $r \in [0, 1-d]$. Let the stratum's score
law under $p$ put mass $r+d$ on a point at or above the cut $c$
and $1-r-d$ below it, and under $q$ put mass $r$ above and $1-r$
below (two-point laws, the cut lying between the two support
points). The laws differ only in the mass at the upper point, so
$\mathrm{TV} = d$; the stratum rates are $r+d$ and $r$, so the
rate difference is exactly $d$. \emph{Sum form.} Apply the
single-stratum construction in both groups with variations
$d_1, d_0$, choosing directions so the two group-rate differences
enter $\mathrm{DP}_p - \mathrm{DP}_q$ with the same sign (group
1's rate higher under $p$ by $d_1$, group 0's lower under $p$ by
$d_0$). Then $|\mathrm{DP}_p - \mathrm{DP}_q| = d_1 + d_0 =
D_{\mathrm{DP}}$: equality in the sum form, the triangle-inequality
step of Theorem~\ref{thm:dp} being tight exactly when the two
signed differences align in this way. \emph{Max form.} Apply the
sum-form equality construction within the label-1 strata only,
with profile $D_{\mathrm{TPR}} \ge D_{\mathrm{FPR}}$, and hold
the label-0 strata identical, so $D_{\mathrm{FPR}} = 0$ and the
FPR-gap difference is $0$. Then $|\mathrm{EO}_p -
\mathrm{EO}_q| = |\mathrm{TPRgap}_p - \mathrm{TPRgap}_q| =
D_{\mathrm{TPR}} = \max(D_{\mathrm{TPR}}, D_{\mathrm{FPR}})$:
equality in the max form. Each stratum needs only rate headroom
$r + d \le 1$, available whenever the profile satisfies $D \le
1$; for $D > 1$ the bound exceeds the maximum possible gap
difference and is vacuous.
\end{proof}

\section{The accounting protocol, for both functionals}
\label{sec:accounting}

An accounting audit asks, per fold of a cross-population design,
whether the realised transport of an audited functional is within
the distributional transport the relevant theorem prices. The protocol was fixed for the FAR in this programme's upgrade cycle (its T4 protocol) and is transported here to the gaps unchanged in form. Fix a functional $G$ (the FAR
of the CTRL forecaster at its FAR-matched threshold; or the EO gap;
or the DP gap) and its transport quantity $D_G$ (the null
step-score TV for the FAR; $D_{\mathrm{EO}}$ / $D_{\mathrm{DP}}$
above). For each leave-one-population-out fold, with test
population the held-out unit and train population the pooled
remainder: compute the realised difference
$r_k = |G_{\mathrm{test}} - G_{\mathrm{train}}|$ and the estimator
$\hat D_{G,k}$ of $D_G$ between the same two populations. The
estimator is the pipelines' histogram TV throughout: score samples
(the conditional samples of Section~\ref{sec:bound} for the gaps;
null step-score samples for the FAR) are clipped to $[0,1]$ and
binned into 50 equal-width bins, and
$\hat{\mathrm{TV}} = \tfrac12 \sum_i |\hat p_i - \hat q_i|$.
A fold \emph{passes} iff $r_k \le \hat D_{G,k} + 0.02$, and the
cell's criterion is \emph{met} iff the median fold
$\hat D_G \ge 0.05$ (the accounting must price non-trivial
transport) and the pass fraction is at least $0.99$. Two honesty
notes are part of the protocol, not footnotes to it. Binning
contracts TV, so $\hat D_G$ estimates $D_G$ from below: a criterion
pass is an empirical fact about the accounting at the registered
slack $0.02$ --- a design constant inherited from the FAR protocol,
not a derived concentration bound --- not a corollary of the
theorems, which bound $r_k$ by the unbinned $D_G$. And a verdict is
a claim about one functional under one theorem: this paper never
lets a FAR verdict stand in for a gap verdict, or vice versa;
Section~\ref{sec:results} attributes each recorded verdict to the
accounting that produced it.

\section{Accounting results}
\label{sec:results}

\paragraph{The FAR accounting, as recorded.} Table~\ref{tab:far}
collects the five recorded verdicts of the earlier cycles. They are
reproduced here as recorded --- each is a verdict about the FAR
functional under the stratum lemma, on the folds named.

\begin{table}[h]
\centering
\begin{tabular}{lllll}
\toprule
Cell & Median TV50 & Folds & Verdict & Fold artifacts \\
\midrule
Family 1, LR & 0.3481 & 50/50 & met & lost (see text) \\
Family 1, GB & 0.1358 & 50/51 & \textbf{not met} & lost (see text) \\
Family 2 2024, LR & 0.4784 & 9/9 & met & survive \\
Family 2 2024, GB & 0.2865 & 9/9 & met & survive \\
Family 2 2025, GB & 0.3475 & 7/7 & met & survive \\
\bottomrule
\end{tabular}
\caption{The FAR accounting (T4 protocol), as recorded in the
programme's cycle reports. The family-1 GB cell missed on one fold
by 0.0034 beyond the slack. The family-1 folds come from an
expanded 51-jurisdiction design whose artifact tree did not survive
an archive migration; the surviving family-1 record is the cycle
reports and an independent tester's recomputation of those cells.}
\label{tab:far}
\end{table}

\paragraph{The gap accounting.} Table~\ref{tab:gap} runs the same
protocol with $G \in \{\mathrm{EO}, \mathrm{DP}\}$ on every fold
partition whose per-record artifacts survive. On family 2 the
criterion is met for both statistics in every cell --- and with
very wide margins, because the realised gap transports there are
tiny (maximum 0.0165): family 2's deployed gaps are small (test-fold
EO gaps at most 0.024) and barely move between forces, so the gap
accounting, while correct, is loose, and the FAR accounting remains
the binding evidence on that programme. On the family-1 folds that
survive --- the six-state leave-one-out partition of the adapted
design's secondary arm, a different and smaller partition than the
lost 51-jurisdiction design --- the criterion is likewise met for
both statistics, and there it does real work: realised EO
transports up to 0.052 are priced by transport quantities with
median 0.110 and minimum fold margin 0.085.

\begin{table}[h]
\centering
\begin{tabular}{lllll}
\toprule
Cell & Statistic & Median $\hat D$ & Pass & Min margin \\
\midrule
Family 2 2024, LR & EO & 0.4149 & 9/9 & 0.2718 \\
 & DP & 0.3135 & 9/9 & 0.2667 \\
Family 2 2024, GB & EO & 0.5011 & 9/9 & 0.3112 \\
 & DP & 0.4427 & 9/9 & 0.2958 \\
Family 2 2025, GB & EO & 0.3442 & 7/7 & 0.2886 \\
 & DP & 0.3241 & 7/7 & 0.2819 \\
Family 1 arm, LR & EO & 0.1104 & 6/6 & 0.0854 \\
 & DP & 0.1020 & 6/6 & 0.0745 \\
\bottomrule
\end{tabular}
\caption{The gap accounting (this paper). Criterion: median
$\hat D \ge 0.05$ and pass fraction $\ge 0.99$; margin
$= \hat D + 0.02 - r_k$ on the tightest fold. Family-2 2025 LR
quantities were also computed (EO median 0.2845, DP 0.2787, 7/7)
but are descriptive only: no recorded verdict exists for that cell,
the LR statistic being degenerate there
(Section~\ref{sec:vacuity}). A 20-bin sensitivity check of the FAR
accounting on family 2 2024 GB failed on one fold (cumbria); the
50-bin convention is the registered one throughout.}
\label{tab:gap}
\end{table}

\paragraph{Re-adjudication.} Applying the rule fixed before the
gap computation --- a recorded verdict follows under the gap bound
iff its own folds' artifacts survive and the gap criterion
reproduces its outcome for both statistics --- the three family-2
verdicts of Table~\ref{tab:far} \textbf{follow}: each is now a
verdict under both accountings. The two family-1 verdicts
\textbf{stand outside} the gap bound, for the archival reason in
Table~\ref{tab:far}: their folds do not exist to be recomputed,
and the six-state computation above, which concerns different
folds, is a corollary about family 1, not a re-adjudication of
them. The family-1 GB \emph{not met} therefore remains what it
always was --- a FAR-accounting fact --- and is not converted,
quietly or otherwise, into a gap statement.

\section{The forecasting programme and its failure}
\label{sec:forecasting}

The programmes above were built to test a stronger claim than
accounting: that a fairness statistic's threshold crossing can be
forecast --- alarmed at least one step ahead, at a controlled
false-alarm rate, with lead times beating trivial baselines. The
claim was pre-registered (FALSIFIER-U) with a binding decision
rule: on the confirmatory arm, the primary forecaster must beat
\emph{all three} baselines (extrapolation, last-value, CUSUM) on
median lead time, with bootstrap intervals excluding ties where the
design requires it. It failed, twice on family 1 and by
non-existence of a test on family 2.

\paragraph{Family 1, absolute threshold.} At the original design
(gap threshold $\tau = 0.05$, horizon $H = 2$ years), TRA-Q's median
lead over covered crossings was 2.0 years --- identical to naive
extrapolation's 2.0 --- at coverage 0.451 of the full-scope
crossing set. The falsifier fired: no lead-time advantage over the
trivial baseline it had to beat.

\paragraph{Family 1, scale-adapted threshold.} A second cycle asked
whether the failure was an artefact of the absolute threshold,
which sits far in the tail of family 2's gap distribution and
arbitrarily within family 1's. Its design adapted the threshold to
the monitored statistic's own scale,
$\tau^* = Q_{0.70}$ of the pooled population gaps --- a constant
anchored to family 1's original design audit (where the operative
0.05 sits at the 69.4th percentile), fixed before any confirmatory
run, and applied to an untouched evaluation. On that arm the test
was fully non-vacuous (7 of 18 trajectories cross after base;
prong-1 and the accounting prong both passed) and the falsifier
fired again, in its most informative form: TRA-Q's median lead was
3.0 years at coverage 0.606, but the margin over extrapolation was
exactly $0.0$ --- a tie, the falsifier's signature --- against
$+1.0$ over last-value (bootstrap CI $[-1.0, 4.0]$, containing 0)
and $+1.0$ over CUSUM (CI $[1.0, 3.0]$, the one criterion met).
The prong required all three; it failed.

\paragraph{Family 2.} Twice, no test existed.
Section~\ref{sec:vacuity} gives the two mechanisms as results; the
facts are these. In 2024, at the headline cell, LR population EO
gaps reach the registered $\tau = 0.05$ in only 2 of 108 force-month
cells, and trajectories cross after the base month only twice
across the threshold grid's first point (then 1, 0, 0): the
scoreability floor of three crossings is not met, for reasons the
crossing-rarity corollary prices exactly. In 2025, the adapted
threshold itself degenerated: $\tau^*(\mathrm{LR}, \mathrm{EO}) =
0.0$ (and likewise for DP), because the deployed LR's scores reach the $0.5$ cut for only $21$ of that year's $65{,}826$ records. In $81$ of the $96$ force-month cells no record reaches the cut at all, so every stratum rate is $0$ and the gaps are identically zero in $84.4\%$ of cells; the remaining $15$ cells hold one to four above-cut records each, moving stratum rates only inside small strata --- gaps at most $0.031$, except one cell whose single above-cut record sits in a five-record stratum and produces a gap of $0.20$. With $84.4\%$ of the pooled cell gaps exactly $0$, the adapted quantile itself is $0$ and there is nothing to forecast. GB in 2025 did produce an evaluable accounting
arm (Table~\ref{tab:gap}'s companion FAR cell) but only one
crossing force: no forecasting adjudication is possible at $n = 1$,
and none is claimed.

\paragraph{Retirement.} The forecasting claim has failed every
non-vacuous test given to it, and the adapted design excludes the
``wrong threshold'' explanation of the first failure. It is
therefore retired, not revised: no further forecaster from this
programme will be proposed on this evidence, and the remainder of
the paper concerns what the programme established instead --- the
bound, the audit, and the two design theorems below.

\section{Two vacuity mechanisms, as theorems}
\label{sec:vacuity}

\begin{corollary}[Crossing rarity]
If a trajectory's per-step crossing probability is at most
$\bar p$ at threshold $\tau$, the probability of any crossing in a
window of $m$ post-base steps is at most $m \bar p$ (union bound),
and a design whose expected crossing count
$(\text{trajectories} \times m \times \bar p)$ falls below its
registered scoreability floor cannot adjudicate forecasting at
$\tau$, whatever forecaster is supplied. Family 2, 2024: the
measured cell share at $\tau = 0.05$ is $\hat p = 2/108 = 0.0185$
over the $9 \times 12$ force-month cells (enumerated cell-wise in
the deposited population table), giving an expected-crossing bound
of $9 \times 11 \times \hat p = 1.83$ against a floor of 3 ---
matching the observed 2 crossings. Family 1's consistency check: the measured cell share at the same threshold is $33/108 = 0.306$ over the $18 \times 6$ task--state--year cells, giving an expected-crossing bound of $18 \times 5 \times 33/108 = 27.5$ over the $18$ trajectories' $5$ post-base steps, and crossings are plentiful --- as observed. In both applications $\hat p$ is the measured share from the same cells to which the corollary is applied --- a plug-in estimate where the corollary is stated for a bound $\bar p$ --- so the two numbers are expected-count diagnostics at the measured share, not strict bounds at a known $\bar p$.
\end{corollary}

\begin{proposition}[Cut degeneracy]
If a deployed rule's scores on a population satisfy
$\Pr[S \ge c \mid Y{=}y, A{=}a] = 0$ for every stratum --- the
scores never reach the decision cut --- then every gap of
Section~\ref{sec:objects} is identically $0$ on that population,
every quantile-adapted threshold of the form
$Q_q(\text{population gaps})$ is $0$ for every $q$ in the range the
adapted design used ($0.60$--$0.80$), and no crossing-based test on
that population has content. Family 2's 2025 LR cell is the instance, cell by cell: in $81$ of the $96$ force-month cells the condition holds exactly --- no record's score reaches the cut, every stratum rate is $0$, and that cell's gaps are identically $0$ --- while the remaining $15$ cells hold the corpus's only $21$ above-cut records and have small nonzero gaps (at most $0.031$, except the one five-record-stratum cell at $0.20$). The pooled conclusion then follows from the zero share, not from the proposition alone: because $81/96 = 84.4\%$ of the pooled cell gaps are exactly $0$, a share exceeding every registered sensitivity quantile, $\tau^* = Q_q(\text{pooled gaps}) = 0$ at $q \in \{0.60, 0.70, 0.80\}$.
\end{proposition}

Both results are elementary; their value is diagnostic. A
monitoring programme can fail to test forecasting for reasons that
are knowable from the design sheet and the score distribution
\emph{before} any forecaster is built, and both of family 2's
failures were of that kind.

\section{The boundary of explanation}
\label{sec:boundary}

Does a bound on the value of history explain the forecasting
failure? The bound in question prices information sets, and must
be stated before it can be evaluated. Write $\Delta_t$ for the
population gap at step $t$, $\mathcal{F}_t$ for the full batch
history to step $t$, and $\mathcal{G}_t =
\sigma(\hat\Delta_{t-1}, \hat\Delta_t)$ for the reference
information: the last two batch gap estimates alone.

\begin{lemma}[History-value bound, at the strength proved]
\label{lem:history}
Let $X$ be any measurable function of the future evolution after
step $t$ (in particular the next $H$ population gaps) and let
$\delta_t := \mathrm{E}\bigl[\mathrm{TV}\bigl(\mathrm{Law}(X
\mid \mathcal{F}_t), \mathrm{Law}(X \mid
\mathcal{G}_t)\bigr)\bigr]$. For any $\mathcal{F}_t$-measurable
alarm rule $A$ (external randomisation allowed), the
$\mathcal{G}_t$-measurable rule $A'$ that fires with probability
$\mathrm{P}(A \text{ fires} \mid \mathcal{G}_t)$ satisfies
$|\mathrm{P}(A \text{ fires}, X \in B) - \mathrm{P}(A' \text{
fires}, X \in B)| \le \delta_t$ for every measurable $B$.
Consequently per-step alarm probabilities, and false-alarm rates
averaged over steps, of the two rules differ by at most $\delta
:= \sup_t \delta_t$; coverage and joint-lead probabilities differ
by at most $(T-1)\delta$ over $T$ monitored steps, and that
clause is carried only at union-bound-sketch level. The stronger
form registered for this lemma --- coverage and joint lead also
at $\delta$ --- is false in general: an independent checker
exhibited disjoint future sub-windows, each at per-step distance
$\delta$, whose coverage probabilities separate by $\approx
k\delta/2$. No verdict in this paper uses the strong form.
\end{lemma}

The proof (a covariance identity, then the variational
definition of TV applied conditionally) and the checker's
obligation-by-obligation adjudication of it are in the deposited
record (Section~\ref{sec:repro}). To size $\delta$ at family 1's
regime, a companion analysis fitted Model G --- a two-channel
Gaussian state-space model of the gap process, $\Delta_{t+1} =
\Delta_t + \mu_t + \eta_t$ with AR(1) level drift $\mu_{t+1} =
\phi\mu_t + \zeta_t$, observed through batches $Y_t = \Delta_t +
\varepsilon_t$ with batch-noise scale $\sigma_n$ --- to the
family's measured increments. Under Model G the future path is
Gaussian under both information sets with a common innovation
covariance, and Pinsker's inequality with Jensen's inequality
gives $\delta \le \sqrt{\bar{\mathrm{KL}}/2}$, where
$\bar{\mathrm{KL}} := \mathrm{E}\bigl[\mathrm{KL}\bigl(
\mathrm{Law}(X \mid \mathcal{F}_t) \,\|\, \mathrm{Law}(X \mid
\mathcal{G}_t)\bigr)\bigr]$ is the mean per-step divergence
between the history-conditional and reference predictive laws:
$\delta$ is the per-step total-variation price of the history
the reference rule cannot see. At the measured regime the
increments' fitted lag-1 and lag-2 autocovariances are
non-positive, which no interior Model G parameterisation fits,
so the drift channel projects to its boundary ($\phi = 0$,
drift variance $0$): the drift channel is exactly zero. The
level channel gives $\bar{\mathrm{KL}} = 0.167$, hence
$\hat\delta = \sqrt{0.167/2} = 0.289$ (the divergence
independently Monte-Carlo-confirmed at $0.167 \pm 0.0003$).
Batch noise is $\sigma_n = 0.037$, about $2.4\times$ the
per-step increment scale $\sigma_{\mathrm{inc}} = 0.015$ ---
history genuinely informs the level, because the noise on each
batch is larger than the movement between batches.

Against that number, the commissioned corollary --- that this
bound explains the two family-1 falsifications --- fails by an
order of magnitude: it would have needed $\delta \le
0.038$--$0.043$ on the covered-only reading and $\le 0.0086$ on
the all-instances reading, and no probed parameterisation of
the fitted model brings $\hat\delta$ below $0.18$. The weakened
coverage clause is no rescue either: at family 1's $T = 6$,
$(T-1)\hat\delta = 1.45$ is vacuous. The bound therefore
\emph{permits} the record rather than explaining it, and the
forecasters still could not beat extrapolation. The two vacuity
events of Section~\ref{sec:vacuity} are covered, but by that
section's own results --- the crossing-rarity corollary's
counting bound and the cut-degeneracy proposition's determinism
argument --- each verified cell-wise against the record by the
same independent checker; the history-value analysis covers
neither falsification, and this paper does not let it borrow
theirs. We report the boundary as a boundary: the explanation
of the forecasting failure is, at present, the failure itself,
documented at pre-registered strength.

\section{Related work}

The monitoring literature's reference system, FairCanary
\cite{faircanary}, detects and attributes bias in a deployed model
after it emerges; it has no forecasting component, which is the gap
this programme set out --- and failed --- to fill. The fairness
statistics are those of \cite{dwork2012,hardt2016}; transport and
drift accounting via score-distribution TV follows the drift
quantities of the FairCanary lineage \cite{faircanary}. Within
this journal, concept-drift work is method-first
\cite{crypsis}: detection-and-mitigation algorithms evaluated on
benchmarks. The present paper differs in kind: it is an audit of a
monitoring programme's own statistics and claims, with the
programme's negative result reported at the same evidentiary
standard as its positive one --- a genre the journal's scope
explicitly houses under evaluating, consolidating, and explaining
discovered knowledge, and under the reporting of successes and
failures of KDD in application.

\section{Limitations}

Four, plainly. (1) The theory is classical: the stratum lemma is
the textbook single-event bound, and the paper claims composition
and audit, not a new inequality. (2) The family-1 expanded-design
artifacts were lost in an archive migration; the affected verdicts
are reported from the cycle record and an independent
recomputation, are attributed to the FAR accounting only, and the
gap re-adjudication of those folds is impossible --- a
data-management failure this paper discloses rather than repairs
retroactively. (3) On family 2 the gap accounting is loose
(Section~\ref{sec:results}); the transport quantities there bound
gap movements an order of magnitude larger than any observed, so
the family-2 gap verdicts certify correctness of the accounting,
not its tightness. (4) EO monitoring presumes labelled batches ---
an audit setting; DP does not --- and both programmes binarise
their protected attributes as their source data require, which
coarsens the fairness questions that can be asked.

\section{Reproducibility}
\label{sec:repro}

Every quantitative claim in this paper traces to a frozen artifact:
the pre-registrations of each cycle (with SHA-256 hashes recorded
at freeze time), the saved per-record score files for both
families, the fold-detail tables of both accountings, and
validator scripts that recompute the headline criteria from those
artifacts (all passing at the time of writing). The theory of
Section~\ref{sec:bound} and the boundary analysis of
Section~\ref{sec:boundary} were additionally checked by an
independent checker working from separately written code, whose
reports are part of the deposited record. The exception is
Limitation (2): the family-1 expanded design's original run files
no longer exist, and its two verdicts rest on the cycle reports
plus the independent tester's recomputation of them. A deposit package, \texttt{fairness-transport-accounting},
has been assembled from these artifacts and validated
end-to-end: the cycles' pre-registrations, the proof records
and both independent checkers' reports, the surviving
fold-detail tables of both accountings, the family-2 population
tables with the per-cell score summary behind
Sections~\ref{sec:forecasting} and~\ref{sec:vacuity}, and a
validator script that recomputes the headline criteria of
Sections~\ref{sec:results} and~\ref{sec:vacuity} from the
deposited tables, all checks passing at the time of writing.
The package is forthcoming as a public repository under the
author's GitHub account
(\texttt{ranadhir19/fairness-transport-accounting}).

\begin{thebibliography}{9}
\bibitem{dwork2012} Dwork, C., Hardt, M., Pitassi, T., Reingold,
O., Zemel, R.: Fairness through awareness. ITCS 2012.
\bibitem{hardt2016} Hardt, M., Price, E., Srebro, N.: Equality of
opportunity in supervised learning. NeurIPS 2016.
\bibitem{faircanary} Ghosh, A., Shanbhag, A., Wilson, C.:
FairCanary: an auditing tool for fair machine learning (AIES
2022; arXiv:2106.07057).
\bibitem{crypsis} Crypsis: an elitism-driven observer-based
approach for detection and mitigation of concept drift. Data
Mining and Knowledge Discovery (2026). DOI
10.1007/s10618-026-01191-1.
\bibitem{bendavid2010} Ben-David, S., Blitzer, J., Crammer, K.,
Kulesza, A., Pereira, F., Vaughan, J.~W.: A theory of learning
from different domains. Machine Learning 79, 151--175 (2010).
\bibitem{policeuk} data.police.uk: UK police stop-and-search open
data, 2024--2025 extracts.
\bibitem{acs} US Census Bureau: American Community Survey,
2014--2019 (folktables construction).
\end{thebibliography}

\end{document}
